\documentclass[11pt, a4paper]{article}

% --- Packages ---
\usepackage[margin=2cm]{geometry} % 2cm margins as per specification
\usepackage{setspace}
\singlespacing % Single line spacing as per specification
\usepackage{graphicx}
\usepackage{hyperref}
\usepackage{titlesec}
\usepackage{caption}

% Set minimum font size for captions/tables to 9pt
\captionsetup{font=small} % 'small' is typically 9pt or 10pt in an 11pt document

\begin{document}

% ==========================================
% TITLE PAGE
% ==========================================
\begin{titlepage}
    \centering
    \vspace*{2cm}

    {\Huge \bfseries Integrating Explainable AI with Deep Learning for Intrusion Detection in IoMT Networks\par}
    \vspace{1.5cm}

    {\Large \bfseries Ralph Lickley\par}
    \vspace{0.5cm}
    {\large Student ID: 730030331\par} % Update if necessary
    \vspace{2cm}

    {\Large \bfseries Abstract\par}
    \vspace{0.5cm}
    \begin{minipage}{0.85\textwidth}
        % SPECIFICATION NOTE: Abstract must be 200-250 words.
        % Briefly summarise the context (IoMT security), the problem (black-box ML models), 
        % the proposed solution (DNN + LIME/SHAP), the methodology (CIC2024 IoMT Dataset), 
        % and your key findings/results.
        [Insert your 200-250 word abstract here.]
    \end{minipage}

    \vfill

    {\Large \bfseries Declaration\par}
    \vspace{0.5cm}
    \begin{minipage}{0.85\textwidth}
        \centering
        \textit{I certify that all material in this dissertation which is not my own work has been identified.}

        \vspace{1cm}
        Signed: \textbf{Ralph James Lickley} \\
        Date: \today
    \end{minipage}

    \vspace{2cm}
\end{titlepage}

% ==========================================
% TABLE OF CONTENTS
% ==========================================
\newpage
\tableofcontents
\newpage

% ==========================================
% BODY OF REPORT (Max 20 Pages from here to Bibliography)
% ==========================================

% --- 1. INTRODUCTION, MOTIVATION, AND AIMS (10%) ---
% SPECIFICATION NOTE: Introduction and background combined must be NO MORE than 2 pages.
\section{Introduction, Motivation, and Aims}

\subsection{Introduction}
The Internet of Medical Things (IoMT) is a subset of the Internet of Things, which is constrained to smart medical equipment
used to monitor patient's health and assist in providing healthcare. For example, devices such as cardiovascular monitors,
temperature trackers, and continuous glucose monitors can all be integrated into the delivery of healthcare. Using these devices
to continuously monitor patients allows healthcare professionals to make more informed decisions for diagnosis and treatment plans.
For example, data derived from IoMT devices has successfully been used to predict the risk people have for heart disease \cite{baseer2024healthcare}.
However, the rapid expansion of IoMT systems has also lead to the attack surface becoming much larger in critical medical
infrastructure. According to the HIPAA Journals' report, during 2025 there were 710 healthcare data breaches, each of
which affecting over 500 people \cite{hipaajournal}. This, alongside the highly sensitive nature of patient's medical data, creates large concerns
over the privacy and handling of patients' data, as well as their safety, given that the equipment used to provide treatment
is now becoming highly interconnected in IoMT. Traditional, signature-based Intrusion Detection Systems (IDS) often fail to detect more
modern, unknown attack signatures, this has lead to Machine Learning (ML) based IDS to becoming a
key area of recent research in network intrusion detection. ML based IDS are able to recognise the complex patterns found in
network traffic to identify malicious activity. However, the lack of interpretability of these ML based
systems limits the widespread adoption, especially in healthcare setting where the safety and privacy of both patients and staff
are critical. To address this lack of interpretability in more complex models, this project focuses on this critical gap by developing
a Deep Neural Network (DNN) IDS, and uses Explainable Artificial Intelligence (XAI) techniques to verify the performance and accuracy of the model.
\subsection{Motivation}
The main motivation behind this project is the requirement for transparency in safety critical systems. Current ML-based
IDS show solid performance, the reasons for decisions they make are still very unclear, this reduces adoption in critical medical
environments because and incorrect classification could lead to disruptions in healthcare delivery, or breaches of patients data.
Existing IDS studies rely on the use of general IoT datasets for training the models, these do not reflect the unique traffic patterns
which are found in IoMT networks, such as the use of lightweight protocols like MQTT. This project uses the CIC2024 IoMT Dataset, this dataset
has been specifically created to benchmark and train IDS for IoMT traffic, it has been generated using traffic collected from real IoMT devices.
Although XAI techniques such as SHAP and LIME have been used in general network security applications, there has been little research applying these
techniques specifically to IoMT scenarios, this project aims to focus on that gap.
\subsection{Background Summary}
% Do not reproduce your literature review verbatim. Briefly summarise:
% 1. IoMT vulnerabilities.
% 2. Machine learning for IDS.
% 3. The role of XAI (LIME and SHAP).


% --- 2. DESIGN, METHODS, AND IMPLEMENTATION (30%) ---
\section{Design, Methods, and Implementation}

\subsection{Dataset Selection and Preprocessing}
% Discuss the CIC2024 IoMT Dataset. Explain your data cleaning, encoding of attack types, 
% and how you handled the massive DDoS/DoS class imbalance (e.g., using SMOTE).

\subsection{Deep Neural Network (DNN) Architecture}
% Detail the design choices of your DNN. 
% Justify your hyperparameter tuning process and success criteria.

\subsection{Integration of Explainable AI (XAI)}
% Explain how you integrated LIME and SHAP into the pipeline. 
% Detail the methods used to extract feature importance.

\subsection{Resource Utilisation and Tools}
% Briefly describe the libraries (e.g., TensorFlow, PyTorch, scikit-learn), hardware, and environments used.


% --- 3. RESULTS, EVALUATION, AND TESTING (30%) ---
\section{Results, Evaluation, and Testing}

\subsection{Experimental Setup and Success Criteria}
% Reiterate the goal (e.g., F1-score >= 0.80) and how testing was conducted.

\subsection{DNN Intrusion Detection Performance}
% Present your classification results. Use legible tables and figures (min 9pt font).
% Discuss F1-score, Precision, Recall, and overall accuracy across the multi-class dataset.
% Compare against baselines (e.g., the Random Forest benchmark from the original CIC2024 paper).

\subsection{XAI Interpretability Validation}
% Present the output from LIME and SHAP. 
% Do they identify the same features for specific attacks? 
% Are the highlighted features logical for the attack type (e.g., DoS attack features)?


% --- 4. REFLECTION, CONCLUSION, AND PRESENTATION (30%) ---
\section{Reflection and Conclusion}

\subsection{Critical Analysis of Results}
% Interpret what your results mean. How successfully did you bridge the gap between high accuracy and interpretability?

\subsection{Project Reflection}
% Fairly examine both the strengths and weaknesses of the project. 
% Discuss project management (referencing your Gantt chart methodology), what went wrong, and how you adapted.

\subsection{Future Work}
% Suggest well-articulated future directions (e.g., testing on real-time hardware, trying different XAI methods).

\subsection{Conclusion}
% Final wrap-up of the project's achievements against the original objectives.


% ==========================================
% BIBLIOGRAPHY
% ==========================================
\newpage
% SPECIFICATION NOTE: Must be at least 11pt font. Will likely take 0.5 to 3 pages.
\bibliographystyle{IEEEtran}
\bibliography{references.bib}

% ==========================================
% APPENDICES (Optional - Does not count towards 20-page limit)
% ==========================================
\newpage
\appendix
\section{Supplementary Diagrams and Code Snippets}
% Use this for extra experimental data, larger diagrams, or system architecture charts 
% that are not strictly necessary for the main narrative but provide completeness.

\end{document}